\documentclass[10pt,a4paper]{amsart}
\usepackage[margin=19mm]{geometry}
\usepackage{amsmath,amssymb,mathtools}
\usepackage[scr=rsfs,cal=euler]{mathalfa}
\usepackage{xfrac}
\usepackage[unicode,hidelinks,bookmarks=false]{hyperref}
\newcommand{\ie}{i.e.\ }
\newcommand{\cf}{cf.\ }
\newcommand{\resp}{respectively\ }
\newcommand{\Subsection}{\subsection}
\providecommand{\nobreakdash}{\nobreak\hbox{-}}
\setlength{\emergencystretch}{2em}
\multlinegap0pt
\usepackage{bm}
\usepackage{bbm}
\usepackage{dsfont}
\usepackage{xspace}
\usepackage{cancel}
\usepackage{mathtools}
\usepackage[shortlabels]{enumitem}
\usepackage{amsthm}
\makeatletter
\@ifundefined{corollary}{\newtheorem{corollary}{Corollary}}{}
\@ifundefined{lemma}{\newtheorem{lemma}{Lemma}}{}
\@ifundefined{proposition}{\newtheorem{proposition}{Proposition}}{}
\@ifundefined{theorem}{\newtheorem{theorem}{Theorem}}{}
\makeatother
\title[Boundedness of solutions in feedback systems with antithetic]{Boundedness of solutions in feedback systems with antithetic controllers}
\author{Moh Kamalul Wafi, Arthur C. B. de Oliveira, Eduardo D. Sontag}
\date{Source: arXiv:2604.27290v2 (CC BY 4.0)}
\begin{document}
\maketitle

\noindent\emph{Source and licence.} Boundedness of solutions in feedback systems with antithetic controllers. Version arXiv:2604.27290v2 (CC BY 4.0), distributed under the Creative Commons Attribution 4.0 International licence, as recorded in the arXiv licence metadata for this exact version and shown on the abstract page https://arxiv.org/abs/2604.27290v2. Full rights notice: licenses/arxiv-2604.27290-CC-BY-4.0.txt.\par\smallskip

\noindent\textit{Editorial scope.} Counted unit: the Corollary whose source label is \texttt{cor:z1\_decrease} in the article (source0.tex lines 685--689, source label \texttt{cor:z1\_decrease}), delivered together with the article's own complete original proof of it (source0.tex lines 691--870, one \texttt{proof} environment of 180 lines). No mathematics is written, repaired or re-derived: statement and proof are the article's own text line for line. Required context retained uncounted: eq:gen\_dynamics (lines 97--107); thm:GenLin-HasProperty (lines 235--238); eq:ca\_dynamics (lines 631--641); prop:T\_fixedpoint (lines 1061--1073); eq:T\_fixedpoint (lines 1063--1065); prop:ell4\_monotone (lines 1129--1138); prop:Lell4\_threshold (lines 1164--1173); lem:p\_above\_theta (lines 1194--1196). Every definition, display and label of those lines is retained unchanged. Omitted from this package and not redistributed: 1--96, 108--234, 239--630, 642--684, 690--690, 871--1060, 1066--1128, 1139--1163, 1174--1193, 1197--1363. Nothing inside a retained excerpt is omitted or altered: every retained body is delivered exactly as the article's lines are written, so no omission receipt of any kind accompanies this package. Citations: the retained excerpts carry no citation command at all, so no bibliography entry of the article is transcribed and no reference is redistributed. Reference adaptation: none was needed and none was made; every reference inside the delivered unit resolves inside it, so no unresolved reference or citation is produced. Printed slips kept literal: no printed word, symbol, hypothesis or number of the counted statement or of its proof was altered. Layout and class: the article's own class is \texttt{article} with packages including PRIMEarxiv, inputenc, fontenc, hyperref, url, booktabs, amsfonts, amsmath, amssymb, mathtools, nicefrac, microtype, lipsum, fancyhdr; the wrapper is the lane's pinned \texttt{amsart} editorial wrapper at 10pt on A4, which restates the theorem-like environments the retained text needs and adds the article's own macro definitions that the retained text needs (none), with extra wrapper packages (bm, bbm, dsfont, xspace, cancel, mathtools, enumitem). The delivered numbering is the unit's own. Assets: no retained span contains a figure environment, a graphics-inclusion command, a PDF-page inclusion, a table or a TikZ picture; no image file is copied into this package, the package declares no asset dependency and no image file is redistributed with it. Licence: version arXiv:2604.27290v2 of this article is distributed under the Creative Commons Attribution 4.0 International licence, as recorded in the arXiv licence metadata for this exact version and shown on the abstract page https://arxiv.org/abs/2604.27290v2. A full rights notice is delivered at licenses/arxiv-2604.27290-CC-BY-4.0.txt. Characters: every retained excerpt is pure ASCII, so the delivered engine, XeLaTeX with its default Latin Modern text font, reports no missing character.
\begin{equation}\label{eq:gen_dynamics}
    \begin{aligned}
    \dot z_1 &= \alpha_1-\alpha_2 z_1 z_2\\[.25em]
    \dot z_2 &= \alpha_3 y-\alpha_4 z_1 z_2
    \end{aligned}
    \qquad \textrm{and} \qquad
    \begin{aligned}
    \dot x &= Ax + bz_1,\\
    y &= c^{\top} x.
    \end{aligned}
\end{equation}

\begin{theorem}
    \label{thm:GenLin-HasProperty}
    For the closed loop system in \eqref{eq:gen_dynamics}, there exist constants $L_0>0$ and $T_0>0$ such that for every $L>L_0$ and $T\geq T_0$, if a solution satisfies $z_1(0)=L$ and $z_1(t)\geq L$ for all $t\in [0,T]$, then $\dot z_1(T)<0$.
\end{theorem}

\begin{equation}\label{eq:ca_dynamics}
    \begin{aligned}
    \dot z_1 &= \alpha_1-\alpha_2 z_1 z_2\\
    \dot z_2 &= \alpha_3 x_2-\alpha_4 z_1 z_2
    \end{aligned}
    \qquad \textrm{and} \qquad
    \begin{aligned}
    \dot x_1 &= \beta_1 z_1-\beta_2 x_1\\
    \dot x_2 &= \beta_3 x_1-\beta_4 x_2,
    \end{aligned}
\end{equation}

\begin{proposition}\label{prop:T_fixedpoint}
    Let $\Delta>0$ and $\alpha_1,\alpha_4>0$. For each $\ell\in\mathbb{R}_+$, consider the fixed-point equation
    \begin{equation}\label{eq:T_fixedpoint}
        \tau(\ell) = \Delta+\frac{\ln(2)}{\alpha_4(\ell+\alpha_1\tau(\ell))}.
    \end{equation}
    Then, $\tau(\ell)$ satisfies:
    \begin{enumerate}
        \item[\emph{(i)}] the unique positive solution $\tau^* = \tau(\ell)$ exists;
        \item[\emph{(ii)}] $\tau(\ell)$ is continuous on $[0,\infty)$, differentiable on $(0,\infty)$, and strictly decreasing; moreover, the right derivative at $\ell=0$ exists and is negative;
        \item[\emph{(iii)}] $\tau(\ell)>\Delta$ for all $\ell\ge0$, and $\lim_{\ell\to\infty}\tau(\ell)=\Delta$.
        \item[\emph{(iv)}] the solution $\tau(\ell;\Delta)$ is strictly increasing in $\Delta$ for each fixed $\ell\ge0$.
    \end{enumerate}
\end{proposition}

\begin{proposition}\label{prop:ell4_monotone}
Let $\tau(\ell):\mathbb{R}_+\to\mathbb{R}_+$ be defined as in 
\eqref{eq:T_fixedpoint}, $\alpha_3>0$ and $G\coloneqq-c^\top A^{-1}b>0$. For each $\ell\in\mathbb{R}_+$, define
\begin{equation*}
    \rho_z(\ell,\tau(\ell))
    := \frac{\alpha_3 G\ell}{4\alpha_4(\ell+\alpha_1\tau(\ell))}.
\end{equation*}
Then $\rho_z:\mathbb{R}_+\to\mathbb{R}_+$ is strictly increasing with 
$\lim_{\ell\to\infty}\rho_z(\ell,\tau(\ell))=\alpha_3G/(4\alpha_4)$.
\end{proposition}

\begin{proposition}\label{prop:Lell4_threshold}
Let $\rho_z(\ell,\tau(\ell))$ be defined as in Proposition~\ref{prop:ell4_monotone}. Then the map 
$\ell\mapsto \ell\rho_z(\ell,\tau(\ell))$ is strictly increasing 
on $[0,\infty)$ and satisfies
\begin{equation*}
    \lim_{\ell\to\infty}\ell\rho_z(\ell,\tau(\ell))=\infty.
\end{equation*}
Consequently, for every $a>0$, there exists a unique $\ell^*>0$ such that 
$\ell^*\rho_z(\ell^*,\tau(\ell^*))=a$.
\end{proposition}

\begin{lemma}\label{lem:p_above_theta}
    Let $\theta:=\alpha_1/\alpha_2$. Consider the dynamics of $z_1(t)$ and $z_2(t)$ in \eqref{eq:gen_dynamics}. Suppose that $z_1(t)\ge L>L_0$ for all $t\in[0,T]$ with $T\ge T_0$. Assume, in addition, that $p(T_0)>\theta$. Then $p(t)>\theta$ for all $t\in[T_0,T]$.
\end{lemma}

\begin{corollary}\label{cor:z1_decrease}
    Consider system~\eqref{eq:ca_dynamics}. There exist constants $L_0>0$ and 
    $T_0>0$ such that for every $L>L_0$ and $T\ge T_0$, if a solution satisfies 
    $z_1(0)=L$ and $z_1(t)\ge L$ for all $t\in[0,T]$, then $\dot z_1(T)<0$.
\end{corollary}

\begin{proof}
The proof follows the same structure as Theorem~\ref{thm:GenLin-HasProperty}, 
but exploits the explicit cascade structure of~\eqref{eq:ca_dynamics}: lower 
bounds on $x_1$, $x_2$, and $z_2$ are derived stage by stage rather than 
via the step-response $g(t)$. We then use $L_0$ to show 
$z_1(T_0)z_2(T_0)>\alpha_1/\alpha_2$ for every $L>L_0$, forcing 
$\dot z_1(T_0)<0$, and verify this domination persists on $[T_0,T]$ while 
$z_1(t)\ge L$ for all $t\in[0,T]$.

Define $\Delta_1 := \ln 2/\beta_2$, $\Delta_2 := \ln 2/\beta_4$, and 
$\delta(\ell,t) := \ln 2/(\alpha_4 U(\ell,t))$ for $\ell\ge0$, $t\ge 0$, 
where $U(\ell,t):=\ell+\alpha_1 t$. For each $\ell\ge0$, define $\tau(\ell)$ 
as the unique positive solution of
\begin{equation*}
    \tau(\ell) = \Delta_1 + \Delta_2 + \delta(\ell,\tau(\ell)),
\end{equation*}
which is guaranteed to exist by Proposition~\ref{prop:T_fixedpoint}, 
taking $\Delta := \Delta_1+\Delta_2$ in that proposition. Then, fix $L_0$ as the unique positive solution of
\begin{equation*}
    \frac{\alpha_3 G\ell^2}{8\alpha_4(\ell+\alpha_1\tau(\ell))} = \frac{\alpha_1}{\alpha_2},
\end{equation*}
whose existence and uniqueness follow from Proposition~\ref{prop:Lell4_threshold},
where $G = \beta_1\beta_3/(\beta_2\beta_4)>0$ is the DC gain.
Once $L_0$ is fixed, we set $T_0:=\tau(L_0)$, which is a fixed positive 
constant. Since $\tau(\ell)$ is strictly decreasing 
(Proposition~\ref{prop:T_fixedpoint}(ii)), for every $L>L_0$ we have 
$\tau(L)<\tau(L_0)=T_0$, so that $T\ge T_0>\tau(L)$ whenever $T\ge T_0$. 
We will next show that the corollary holds for these choices of $L_0$ and $T_0$.

Let $L>L_0$, $T\ge T_0$ and suppose
\begin{equation}\label{eq:z1_low}
    z_1(t)\ge L, \qquad \text{for all } t\in[0,T].
\end{equation}
We derive lower bounds for $x_1(t)$, $x_2(t)$, and $z_2(t)$, each holding 
after a finite delay within $[0,T_0]\subset[0,T]$.

\begin{enumerate}[leftmargin=*]
    \item \textbf{Estimate of $x_1(t)$}. On $[0,T]$, using 
    $\dot x_1=\beta_1 z_1-\beta_2 x_1$ together with~\eqref{eq:z1_low}, 
    we obtain
    \begin{equation*}
        \dot x_1 \ge \beta_1 L - \beta_2 x_1.
    \end{equation*}
    Define $\dot w_1(t)=\beta_1 L-\beta_2 w_1(t)$ with $w_1(0)=x_1(0)$. 
    By the comparison principle, this leads to $x_1(t)\ge w_1(t)$ for all
    $t\in[0,T]$. Hence
    \begin{equation*}
        x_1(t)\ge \frac{\beta_1 L}{\beta_2}\Bigl[1-e^{-\beta_2 t}\Bigr] 
        + x_1(0)e^{-\beta_2 t}
        \ge \frac{\beta_1 L}{\beta_2}\Bigl[1-e^{-\beta_2 t}\Bigr].
    \end{equation*}
    Since $1-e^{-\beta_2 t}$ is increasing in $t$. By defining $\Delta_1\coloneqq(\ln 2)/\beta_2$, this yields
    \begin{equation}\label{eq:x1_low}
        x_1(t)\ge \frac{\beta_1 L}{\beta_2}\Bigl[1-e^{-\beta_2 t}\Bigr]\ge
        \frac{\beta_1 L}{\beta_2}\Bigl[1-e^{-\beta_2\Delta_1}\Bigr] 
        = \frac{\beta_1 L}{2\beta_2} \eqqcolon \bar\rho_1(L)
        \qquad\text{for all } t\in[\Delta_1,T].
    \end{equation}

    \item \textbf{Estimate of $x_2(t)$}. On the interval $[\Delta_1,T]$, 
    using $\dot x_2=\beta_3 x_1-\beta_4 x_2$ and~\eqref{eq:x1_low} 
    results in
    \begin{equation*}
        \dot x_2 \ge \beta_3\bar\rho_1(L) - \beta_4 x_2.
    \end{equation*}
    Define $\dot w_2(t)=\beta_3\bar\rho_1(L)-\beta_4 w_2(t)$ with 
    $w_2(\Delta_1)=x_2(\Delta_1)$. By the comparison principle, we have 
    $x_2(t)\ge w_2(t)$ for all $t\in[\Delta_1,T]$. Hence
    \begin{equation*}
        x_2(t) \ge \frac{\beta_3}{\beta_4}\bar\rho_1(L)
        \Bigl[1-e^{-\beta_4(t-\Delta_1)}\Bigr]
        + x_2(\Delta_1)e^{-\beta_4(t-\Delta_1)}
        \ge \frac{\beta_3}{\beta_4}\bar\rho_1(L)
        \Bigl[1-e^{-\beta_4(t-\Delta_1)}\Bigr].
    \end{equation*}
    Since $t-\Delta_1\ge0$, thus $1-e^{-\beta_4(t-\Delta_1)}$ is increasing in $t$. Defining $\Delta_2\coloneqq(\ln 2)/\beta_4$, this results in
    \begin{equation}\label{eq:x2_low}
        x_2(t)\ge\frac{\beta_3}{\beta_4}\bar\rho_1(L)
        \Bigl[1-e^{-\beta_4(t-\Delta_1)}\Bigr] \ge 
        \frac{\beta_3}{\beta_4}\bar\rho_1(L)
        \Bigl[1-e^{-\beta_4\Delta_2}\Bigr]
        =\frac{\beta_3}{2\beta_4}\bar\rho_1(L)
        \eqqcolon \bar\rho_2(L) 
        \qquad\text{for all } t\in[\Delta_1+\Delta_2,T].
    \end{equation}

    \item \textbf{Upper bound for $z_1$.}
    Before lower bounding $z_2$, we must first derive an upper bound for 
    $z_1$ in the interval. From the equation $\dot z_1=\alpha_1-\alpha_2 
    z_1 z_2$ and the nonnegativity $z_1,z_2\ge 0$, we obtain
    $\dot z_1(t)=\alpha_1-\alpha_2 z_1 z_2\le\alpha_1$ for all $t\ge 0.$
    Integrating on any interval $[0,t]$ yields
    \begin{equation*}
        z_1(t) \le z_1(0) + \alpha_1 t.
    \end{equation*}
    Using the assumption $z_1(t)\ge L$ for all $t\in[0,T_0]$ and 
    $z_1(0)=L$, then for every $t\in[0,T_0]$,
    \begin{equation}\label{eq:z1_window}
        z_1(t)\le L+\alpha_1 t
        \le L+\alpha_1T_0 = U(L,T_0).
    \end{equation}
    Thus, on any time interval of length $T_0$ on which $z_1$ remains 
    above $L$, it is also bounded from above by $U(L,T_0)$.

    \item \textbf{Estimate of $z_2(t)$}. On the interval 
    $[\Delta_1+\Delta_2,T_0]$, using $\dot z_2=\alpha_3 x_2-\alpha_4 
    z_1 z_2$, \eqref{eq:x2_low}, and~\eqref{eq:z1_window} yield
    \begin{equation*}
        \dot z_2(t) \ge \alpha_3\bar\rho_2(L) - \alpha_4 U(L,T_0)\,z_2(t).
    \end{equation*}
    Define $\dot v(t)=\alpha_3\bar\rho_2(L)-\alpha_4 U(L,T_0)v(t)$ with 
    $v(\Delta_1+\Delta_2)=z_2(\Delta_1+\Delta_2)$.
    By the comparison principle we have $z_2(t)\ge v(t)$ for all
    $t\in[\Delta_1+\Delta_2,T_0]$. Hence, with $\bar U\coloneqq U(L,T_0)$,
    \begin{equation*}
        \begin{aligned}
            z_2(t) &\ge \frac{\alpha_3}{\alpha_4\bar U}\bar\rho_2(L) 
            \Bigl[1-e^{-\alpha_4\bar U(t-(\Delta_1+\Delta_2))}\Bigr]
            + z_2(\Delta_1+\Delta_2)\,e^{-\alpha_4\bar U(t-(\Delta_1+\Delta_2))} \\
            &\ge \frac{\alpha_3}{\alpha_4\bar U}\bar\rho_2(L) 
            \Bigl[1-e^{-\alpha_4\bar U(t-(\Delta_1+\Delta_2))}\Bigr].
        \end{aligned}
    \end{equation*}
    Since $t-(\Delta_1+\Delta_2)\ge 0$, thus $1-e^{-\alpha_4\bar U(t-(\Delta_1+\Delta_2))}$ is increasing in $t$. Using $\delta(L,T_0)=(\ln 2)/(\alpha_4 U(L,T_0))$,
    \begin{equation}\label{eq:z2_low_4d}
        z_2(t) \ge \frac{\alpha_3}{2\alpha_4 U(L,T_0)}\bar\rho_2(L) 
        = \frac{\alpha_3 GL}{8\alpha_4(L+\alpha_1 T_0)} 
        \eqqcolon \bar\rho_z(L,T_0)
        \qquad\text{for all } t\in[\Delta_1+\Delta_2+\delta(L,T_0),T_0].
    \end{equation}
\end{enumerate}

Note that the bounds on $x_1$ and $x_2$ persist while $z_1(t)\ge L$, and 
hence are valid on $[\Delta_1,T]$ and $[\Delta_1+\Delta_2,T]$ respectively. 
For $z_2$, however, the damping coefficient $\alpha_4 z_1$ depends on $z_1$ 
itself, so we need the upper bound $z_1(t)\le U(L,T_0)$, which is only 
available on $[0,T_0]$. This is why the estimate for $z_2$ is restricted 
to $[0,T_0]$.

Unlike the general case, where $y\ge GL/2$ follows from $g(t)\ge G/2$, 
here the bounds on $x_1$, $x_2$, and $z_2$ are derived explicitly stage 
by stage through the cascade $z_1\to x_1\to x_2\to z_2$. 
From~\eqref{eq:z1_low} and~\eqref{eq:z2_low_4d}, 
$z_1(T_0)z_2(T_0)\ge L\bar\rho_z(L,T_0)$.
By the choice of $L_0$ and since $\ell\mapsto\ell\bar\rho_z(\ell,T_0)$ 
is strictly increasing in $\ell$, for every 
$L>L_0$,
\begin{equation*}
    L\bar\rho_z(L,T_0) > L_0\bar\rho_z(L_0,T_0) = \frac{\alpha_1}{\alpha_2},
\end{equation*}
where the last equality holds by definition of $L_0$. Thus
$z_1(T_0)z_2(T_0)>\alpha_1/\alpha_2$ and 
$\dot z_1(T_0)=\alpha_1-\alpha_2 z_1(T_0)z_2(T_0)<0$.

However, as $T$ increases beyond $T_0$, the bound $\bar\rho_z(L,T)$
decreases in $T$, so this estimate does not directly extend beyond $T_0$.
As in the proof of Theorem~\ref{thm:GenLin-HasProperty}, introduce the product
$p(t):=z_1(t)z_2(t)$ and $\theta\coloneqq\alpha_1/\alpha_2$.
We claim that $p(t)>\theta$ for all $t\in[T_0,T]$.
Indeed, the proof is identical to that of Lemma~\ref{lem:p_above_theta}, with
$\rho_z$ replaced by $\bar\rho_z$ and
$L^*:=2\alpha_4\theta/(\alpha_3G)$ replaced by
$L_*:=4\alpha_4\theta/(\alpha_3G)$.
    \footnote{The factor-of-two difference 
    between $L^*$ and $L_*$ reflects the difference in the lower bound on $y$: 
    in the general case, $y\ge GL/2$ follows from the step-response estimate 
    $g(t)\ge G/2$, while here the explicit cascade gives 
    $x_2\ge\bar\rho_2(L)=GL/4$, 
    which carries an additional factor of $1/2$ through the two cascade stages. 
    The remainder of the argument is unchanged: one shows 
    $L_*\bar\rho_{z}(L_*,\tau(L_*))<\theta$, which implies $L_0>L_*$, and 
    hence $\alpha_3\bar\rho_2(L_0)-\alpha_4\theta>0$, giving 
    $\dot p(t_0)>0$ and contradicting $\dot p(t_0)\le 0$.}
Therefore,
\begin{equation*}
    \dot z_1(t) = \alpha_1-\alpha_2 p(t) < \alpha_1-\alpha_2\theta = 0,
    \qquad\text{for all } t\in[T_0,T].
\end{equation*}
In particular, $\dot z_1(T)<0$, completing the proof.
\end{proof}

\end{document}
